\chapter{Configuration and Development Record}
\label{app:devrecord}

This appendix holds the firmware configuration, the record of development faults, and supporting detail moved out of Chapter~\ref{chap:design}. The full text of the earlier, chronological Methodology is kept in the project repository (\texttt{Report/\allowbreak{}archive\_\allowbreak{}full\_\allowbreak{}2026-10-11/}). Figure~\ref{fig:timeline} gives the project timeline as run, for the phases of Section~\ref{sec:phases}.

\begin{figure}[htbp]
\centering
% Project timeline as run, 3 August to 21 October 2026. One unit on x is one day from 3 August.
% Sources: git log (commit dates), PROJECT_LOG.md section 3, CHANGELOG.md, Planning/lab-notebook.md,
% Report/appendixc.tex (evidence index dates). Bars span the first and last dated record of each activity.
\begin{tikzpicture}[x=0.12cm, y=1cm]
  \def\rowh{0.62}
  % ---- week grid and labels (weeks start on Monday) ----
  \foreach \d/\day/\mon in {0/3/Aug, 7/10/Aug, 14/17/Aug, 21/24/Aug, 28/31/Aug, 35/7/Sep, 42/14/Sep, 49/21/Sep, 56/28/Sep, 63/5/Oct, 70/12/Oct, 77/19/Oct} {
    \draw[swGrey!40, line width=0.3pt] (\d,0.35) -- (\d,-9.0);
    \node[font=\swfont, text=swInk, anchor=south, align=center, inner sep=1pt] at (\d+3.5,0.35) {\day\\\mon};
  }
  \draw[swGrey!40, line width=0.3pt] (80,0.35) -- (80,-9.0);
  % ---- rows: label, then bars (start day, end day) ----
  \newcommand{\swrow}[2]{\node[font=\swfont, text=swInk, anchor=east, align=right] at (-1,-#1*\rowh) {#2};}
  \newcommand{\swbar}[4][swBlue]{\fill[#1!35, draw=#1, line width=0.4pt] (#3,-#2*\rowh-0.17) rectangle (#4+1,-#2*\rowh+0.17);}
  \newcommand{\swmile}[3]{\node[draw=swAmber, fill=swAmberFill, diamond, inner sep=1.6pt] at (#2+0.5,-#1*\rowh) {};
    \node[font=\swfont, text=swInk, anchor=west, inner sep=1pt] at (#2+1.6,-#1*\rowh) {#3};}

  \swrow{0}{1. Literature review and research search}
  \swbar{0}{3}{12}  % papers 6 to 15 Aug
  \swbar{0}{47}{47} % automated search and gap ledger, 19 Sep
  \swbar{0}{65}{65} % review aligned with the dashboard scope, 7 Oct

  \swrow{1}{Planning, components, vehicle access}
  \swbar{1}{28}{36} % 31 Aug to 8 Sep (PROJECT_LOG)

  \swrow{2}{2. Test procedures and criteria}
  \swbar{2}{57}{58} % TEST_PROCEDURES.md, 29 to 30 Sep
  \swbar{2}{69}{69} % T15 added, 11 Oct
  \swbar{2}{73}{73} % C11 added, 15 Oct

  \swrow{3}{3. Telemetry firmware}
  \swbar{3}{24}{24} % baseline 27 Aug
  \swbar{3}{44}{79} % dual-IMU FreeRTOS 16 Sep to v0.7.4.1 21 Oct

  \swrow{4}{3. Camera units, detector, ToF fusion}
  \swbar{4}{44}{44} % MVP 16 Sep
  \swbar{4}{56}{73} % first live run 28 Sep to time to contact 15 Oct

  \swrow{5}{3. Pothole model (Colab)}
  \swbar{5}{60}{69} % runs 2 and 6 Oct, C10 11 Oct

  \swrow{6}{3. Dashboard integration}
  \swbar{6}{56}{73} % camera page 28 Sep to two-sided Surroundings 15 Oct

  \swrow{7}{3. Enclosures and mounts}
  \swbar{7}{47}{47} % OpenSCAD camera mount 19 Sep
  \swbar{7}{70}{77} % SolidWorks 12 Oct to printing on 19 Oct

  \swrow{8}{4. Start-up checks and fault diagnosis}
  \swbar[swGreen]{8}{53}{74} % boot loop 25 Sep to housing recurrence 16 Oct

  \swrow{9}{5. Bench integration and acceptance tests}
  \swbar[swGreen]{9}{53}{73} % T1 25 Sep to Pi 5 detector 15 Oct
  \swbar[swGreen]{9}{79}{79} % wired UART check 21 Oct

  \swrow{10}{6. Vehicle work}
  \swbar[swAmber]{10}{58}{59} % BMS frames on the vehicle, 30 Sep to 1 Oct
  \swbar[swAmber]{10}{77}{79} % mounting, road session 19 Oct, data check 20 Oct, 21 Oct

  \swrow{11}{Analysis and report}
  \swbar[swGrey]{11}{4}{4}   % Introduction started 7 Aug
  \swbar[swGrey]{11}{53}{79} % report written with the work from 25 Sep

  % ---- milestones, one row each so that the labels do not overlap ----
  \swrow{12}{Milestones}
  \swmile{12}{59}{}
  \node[font=\swfont, text=swInk, anchor=east, inner sep=1pt] at (58,-12*\rowh) {1 Oct: dashboard replaces the HUD};
  \swmile{13}{73}{}
  \node[font=\swfont, text=swInk, anchor=east, inner sep=1pt] at (72,-13*\rowh) {15 Oct: two-sided layout};
  \swmile{14}{77}{}
  \node[font=\swfont, text=swInk, anchor=east, inner sep=1pt] at (76,-14*\rowh) {19 Oct: first road session};
\end{tikzpicture}

\caption{Project timeline as run, 3 August to 21 October 2026, by week. Each bar spans the first and last dated record of an activity in the git history, \texttt{PROJECT\_LOG.md}, \texttt{CHANGELOG.md}, \texttt{Planning/lab-notebook.md} and the evidence index (Appendix~\ref{app:evidence}); short bars are single dated events. The numbers are the phases of Figure~\ref{fig:vmodel}. Blue: definition and building; green: bench checks and tests; amber: vehicle work; grey: analysis and writing.}
\label{fig:timeline}
\end{figure}

\section{Firmware configuration}

Table~\ref{tab:config} gives the telemetry unit's configuration as committed for firmware v0.7.4, flashed on 11 October 2026 and used for test T15 and the first road session; Table~\ref{tab:fw-versions} gives the later changes. It is taken from \texttt{TelemetryUnit/\allowbreak{}platformio.ini} and \texttt{TelemetryUnit/\allowbreak{}src/\allowbreak{}main.cpp}; installed package versions are those in the PlatformIO folders on the development laptop, since \texttt{platformio.ini} does not pin them. Against v0.7.3, used for the bench results before 11 October, v0.7.4 changes the IMU low-pass filter from 260 to 94\,Hz and adds the record identity and the store-and-forward ledger (rows marked v0.7.4); tasks and queues are unchanged.

\begin{longtable}{|>{\raggedright\arraybackslash}p{4.0cm}|>{\raggedright\arraybackslash}p{9.8cm}|}
\caption{Telemetry unit configuration (firmware v0.7.4).}
\label{tab:config}\\
\hline
\textbf{Item} & \textbf{Setting} \\ \hline
\endfirsthead
\hline
\textbf{Item} & \textbf{Setting} \\ \hline
\endhead
Build environment & PlatformIO, environment \texttt{esp32-s3-devkitc-1}, Arduino framework \\ \hline
Platform and core (installed) & \texttt{espressif32} platform 7.0.1; \texttt{framework-arduinoespressif32} 3.20017.241212 (Arduino-ESP32 core 2.0.17) \\ \hline
Libraries (installed) & Adafruit MPU6050 2.2.9 (declared \^{}2.2.4); Adafruit Unified Sensor 1.1.15 (\^{}1.1.14); Adafruit BusIO 1.17.4 (\^{}1.14.1) \\ \hline
Build flags & \texttt{ARDUINO\_USB\_MODE=1}, \texttt{ARDUINO\_\allowbreak{}USB\_\allowbreak{}CDC\_\allowbreak{}ON\_\allowbreak{}BOOT=1} (application serial on native USB) \\ \hline
Flash & 16\,MB, DIO mode; PSRAM memory type left at default (PSRAM not used) \\ \hline
Partition table & \texttt{default\_16MB.csv}: nvs 0x9000 (20\,KB), otadata 0xE000 (8\,KB), app0 0x10000 and app1 0x650000 (6.25\,MB each), data 0xC90000 (3.375\,MiB, LittleFS), coredump 0xFF0000 (64\,KB) \\ \hline
Serial & USB CDC at 115200\,bit/s; bounded 3\,s wait for a host at boot \\ \hline
IMU1 & MPU6050 at 0x68, I\textsuperscript{2}C bus 0, SDA GPIO8, SCL GPIO9 \\ \hline
IMU2 & MPU6050 at 0x68, I\textsuperscript{2}C bus 1, SDA GPIO17, SCL GPIO18 \\ \hline
I\textsuperscript{2}C clock & 400\,kHz on both buses \\ \hline
IMU settings & $\pm 8\,g$, $\pm 500^{\circ}$/s, low-pass filter 94\,Hz (v0.7.4; 260\,Hz up to v0.7.3); 200\,Hz (5\,ms) sampling \\ \hline
IMU sample record & \texttt{sensor\_id} (uint8, 0 or 1); \texttt{sequence} (uint32, per sensor); \texttt{timestamp\_us} (int64, \texttt{esp\_timer\_get\_time()}); \texttt{ax, ay, az} (float, m/s\textsuperscript{2}); \texttt{gx, gy, gz} (float, rad/s) \\ \hline
Calibration & 200 samples at rest; factor accepted in 0.5 to 2.0; stored in NVS namespace \texttt{imucal} \\ \hline
Modem & A7670X on UART1, RX GPIO47, TX GPIO48, 115200\,bit/s 8N1; PWRKEY GPIO4, RESET GPIO5 \\ \hline
Sync pulse & GPIO15, toggled every 500\,ms \\ \hline
Hall counter & Pulse counter unit 0, both edges, filter 1000 APB cycles (12.5\,$\mu$s); input on GPIO15 (loopback) until the tap is wired; pole pairs not set \\ \hline
BMS & BLE, service 0xFF00, notify 0xFF01, write 0xFF02; query every 3\,s; receive buffer 128 bytes \\ \hline
Periods & GNSS poll 3\,s; statistics 2\,s; ride window 1\,s (200 samples); dashboard line 1\,s (at most 480 bytes); MQTT publish 10\,s; network position 60\,s without a fix; assisted GNSS every 2\,h \\ \hline
Fusion & Pair skew at most 0.5\,s; fallback after 3\,s; $\sigma$ floor 0.05\,m/s\textsuperscript{2}; speed index $I_v$ needs $v \ge 1$ \\ \hline
MQTT & Broker \texttt{test.mosquitto.org}:1883, topic \texttt{sw7/telemetry}, QoS~1, keep-alive 60\,s; retry 30\,s doubling to 300\,s; session lost after 3 failed publishes \\ \hline
Local log & LittleFS \texttt{/telemetry.log}; cap 3\,MB; 8\,KB memory buffer flushed every 60\,s \\ \hline
Record identity (v0.7.4) & \texttt{dev}, \texttt{boot}, \texttt{seq} on every record, \texttt{utc\_ms} from GNSS once fixed; replays marked \texttt{"replay":true} \\ \hline
Store-and-forward (v0.7.4) & NVS namespace \texttt{sw7sync} (boot, seqblk, acked, lost); sequence numbers reserved in blocks of 64; replay at most one record per 300\,ms, 25 log lines per pass, 5\,s after a failed replay; unacknowledged records removed by the log cap counted as lost \\ \hline
\multicolumn{2}{|l|}{Tasks (name: core, priority, stack in bytes)} \\ \hline
Acquisition & Imu1Task, Imu2Task: core 1, priority 3, 4096 \\ \hline
Statistics, sync & StatsTask: core 1, priority 2, 4096; SyncTask: core 1, priority 2, 2048 \\ \hline
Ride features & RideChar1Task, RideChar2Task: core 1, priority 1, 4096 \\ \hline
Modem, BMS, dashboard, Hall & ModemTask: core 0, priority 1, 4096; BmsTask: core 0, priority 1, 8192; DashTask: core 0, priority 1, 4096; HallTask: core 0, priority 1, 2048 \\ \hline
\multicolumn{2}{|l|}{Queues (depth in items)} \\ \hline
Queues & imu1Queue, imu2Queue: 100 samples each; imu1CharQueue, imu2CharQueue: 400 each; gnssQueue: 20; bmsQueue: 20 \\ \hline
\end{longtable}

Table~\ref{tab:fw-versions} gives the changes after v0.7.4, from \texttt{TelemetryUnit/\allowbreak{}CHANGELOG.md} and the lab notebook.

{\small
\begin{longtable}{|>{\raggedright\arraybackslash}p{1.5cm}|>{\raggedright\arraybackslash}p{8.6cm}|>{\raggedright\arraybackslash}p{3.6cm}|}
\caption{Firmware changes after v0.7.4 and their status on 22 October 2026.}
\label{tab:fw-versions}\\
\hline
\textbf{Version} & \textbf{Change} & \textbf{Status} \\ \hline
\endfirsthead
\hline
\textbf{Version} & \textbf{Change} & \textbf{Status} \\ \hline
\endhead
v0.7.4.1 & Every dashboard line also on UART2 (IO41 TX, IO42 RX, 460\,800\,baud); a line is dropped and counted if the 2\,KB buffer is full; local commands accepted from the Pi & Flashed 21 Oct; 2,489 lines, 0 bad \\ \hline
v0.7.5 & Raw 200\,Hz stream (\texttt{RAW}, \texttt{RAWH}, \texttt{RAWS} lines with CRC-16); read-failure and stuck counters; \texttt{front\_node} build (one IMU, no modem, BLE or GNSS); no position in the published record & Front node flashed 22 Oct; host tests 2,326 raw-stream checks \\ \hline
v0.7.6 & Statistics and ride tasks on core~0 at priority~2; IMU catch-up bounded at 20 periods (late count); rates over real elapsed time; dashboard fields \texttt{cons\_age\_ms}, \texttt{qmax}, \texttt{late}, \texttt{idle}; \texttt{!modem off} and \texttt{!modem on}; front sync on GPIO5 (build flag) & Both nodes flashed 22 Oct; host test 254 pacing checks \\ \hline
v0.7.7 & Automatic IMU re-initialisation after 1\,s of failed or stuck reads, back-off 1\,s then 5\,s; I\textsuperscript{2}C timeout 10\,ms and bus restart; fields \texttt{reinit}, \texttt{imu\_up}; IMU2 off in the rear build & Built, host-tested; flashed to both nodes 22 Oct \\ \hline
v0.7.8 & Command task: commands not using the modem, MQTT, flash log or ledger run at once, the rest queued for the modem task, which cuts its long waits short; field \texttt{cmd} & Built, host-tested; flashed to both nodes 22 Oct \\ \hline
\end{longtable}
}

\section{Development faults}

Table~\ref{tab:faults} lists every fault found during development, with its cause and fix. Eight of them, which needed a diagnosis, are described step by step in Section~\ref{sec:investigations}.

\begin{longtable}{|>{\raggedright\arraybackslash}p{3.6cm}|>{\raggedright\arraybackslash}p{5.0cm}|>{\raggedright\arraybackslash}p{4.9cm}|}
\caption{Development faults, causes and fixes.}
\label{tab:faults}\\
\hline
\textbf{Symptom} & \textbf{Cause (how found)} & \textbf{Fix and lesson} \\ \hline
\endfirsthead
\hline
\textbf{Symptom} & \textbf{Cause (how found)} & \textbf{Fix and lesson} \\ \hline
\endhead
Boot loop at a fixed program counter after enabling 16\,MB flash and PSRAM & Quad-I/O access requested on flash that runs in DIO mode (one setting changed at a time; ROM banner reports DIO; chip detection confirmed 8\,MB PSRAM) & 16\,MB with DIO mode set together with the partition table; PSRAM setting left at default. Set flash size, mode and table together \\ \hline
No output past the ROM banner on the target board & Serial had been routed to UART0 for the bare development board; on the target board UART0 goes to a separate CH340K port (each hypothesis tested in turn; schematic) & Native-USB flags restored and documented in \texttt{platformio.ini}. The two boards need opposite settings \\ \hline
Modem silent after a reset & Power key not pulsed (v0.1); later, the pulse at every boot switched an already running modem off & v0.2.0 pulses GPIO4 for 3\,s; v0.5.2 pulses only if \texttt{AT} gets no answer \\ \hline
GNSS queue full after about a minute & Statistics task drained the IMU queues but not the GNSS queue & GNSS queue added to the drain cycle \\ \hline
Every GNSS poll rejected, then the first fix at 0.25 and 0.10 degrees & Assumed 13-field reply; real reply 9 fields indoors, 17 or 18 with a fix, decimal degrees (raw replies logged) & Parser anchored on the hemisphere letter, layout and jump checks (v0.5.1 to v0.5.3) \\ \hline
Zero GNSS polls reported for the first minutes & Polls time out for 60 to 90\,s after GNSS power-on; statistics count only answered polls & Timeout lines distinguish the cases; allow warm-up before recording \\ \hline
Sync output pin not on any header & Placeholder chosen before the schematic was available (schematic and inspection) & Reassigned to GPIO15 before wiring \\ \hline
No sync edges, later a burst of spurious edges & Jumper disturbed by handling (ESP32 log showed the pulse still running) & Reseated; secured connector needed on the vehicle; off-pattern edges discarded \\ \hline
BMS scan never retried & Scan made conditional on an address set only after the device is found & Scan whenever disconnected and not scanning \\ \hline
BMS connection dropped before reading & Characteristics not under the assumed service in this BLE library (compared with a teammate's parser) & All discovered services searched \\ \hline
Every BMS frame rejected & Checksum included the register byte, so it failed by exactly 3 (the firmware's own query frame shows the rule) & v0.5.4; confirmed by 14 valid records on 30 Sep \\ \hline
MQTT reported connected when it was not; no retry after a failure & First \texttt{OK} taken as success; separate result line ignored (modem manual) & v0.5.0 waits for the result code; retry with backoff \\ \hline
Reconnection 36\,s after coverage returned & Waited for the next scheduled retry although data were back in 2.4\,s & v0.5.3 reconnects on the modem's data-available report (about 5\,s) \\ \hline
33 to 40\,ms acquisition gaps every 10\,s & Flash writes suspend code running from flash on both cores (logging on and off test) & v0.6.0 buffers and flushes once a minute; pause length unchanged \\ \hline
Most stored records not valid JSON & Missing values printed as \texttt{nan} (full log read-back) & v0.6.3 writes \texttt{null} \\ \hline
Speed index dominated by a constant & Mean square of the magnitude includes gravity & v0.7.2 uses the variance \\ \hline
A person reported as ``motorcycle''; poor confidence & Placeholder label shifted every class; camera buffers in BGR order & Label list and channel order corrected (fixed together) \\ \hline
16 alert activations in 47\,s for one person & Alert cleared on a single missed frame near the threshold & Clears after 5 frames; nearer band needs 3 frames \\ \hline
Detector slower with the live view on & Drawing and encoding in the detection loop & Moved to a separate thread taking the newest frame \\ \hline
Pink image; later black objects purple & Module has no infrared-cut filter (stills with standard and NoIR tuning) & NoIR tuning and saturation 1.8; dark-pixel correction in the live view only \\ \hline
Raspberry Pi~4 drive would not boot & Boot files copied instead of an image written: one partition, missing root & Rewritten with Raspberry Pi Imager \\ \hline
Raspberry Pi~4 did not join the hotspot; ssh closed & 5\,GHz hotspot without a regulatory domain; ssh not re-enabled in a second provisioning pass & Hotspot on 2.4\,GHz and country code set (sufficient fix not isolated); ssh enabled \\ \hline
Recording aimed at the ceiling & No rigid camera mount & Live preview before recording; printed casing \\ \hline
Bracket model with no lens lip & Cavity cut through the face, visible only in a full render & Cut bounded; checked by full render (bracket not printed) \\ \hline
Simulated speed and battery on the dashboard with no sign & Dashboard fell back to simulation when the bridge address changed & Live-only mode and network discovery \\ \hline
Pi~4 not found on a phone hotspot; 6,828 lookup warnings & 1\,s probe shorter than the hotspot round trip; searches every 2\,s & 3\,s probe; back-off to 30\,s; beacon tried at once \\ \hline
Dashboard crash about 2\,s after showing the web map & Native crash in the web engine on the Pi~5 & Native map selected, web map not created; Qt tile map \\ \hline
Screen blank after 600\,s; poller at 53\% of a core & X11 screensaver and power management; game-controller poller & Both switched off by the start script \\ \hline
Team's dashboard started after a reboot & Desktop auto-start is not read on the bare X display & System service for the integrated copy \\ \hline
All three distances null & Start-up stopped at the first sensor that did not answer & Skip, log and re-probe; 0.5\,s I/O timeout \\ \hline
Pi~4 under-voltage, throttling, failure to boot & LTE modem powered from the Pi's USB port (2\,s power-flag log) & Separate supply for the telemetry unit; fan for temperature \\ \hline
Touch felt slow & Panel reports contact with the first position; buttons act on release (25 taps read directly) & Act on press and hide the pointer (next step) \\ \hline
Rear drop counts rising at about 168 and 46 a second, \texttt{imu\_hz} frozen at 200, vibration null & Statistics and ride tasks starved on the IMU core (Section~\ref{sec:inv-stall}) & v0.7.6 moves them to core~0; health fields in the dashboard line \\ \hline
IMU repeats one reading after its wires are refitted & Sensor back in its power-on sleep state; configured only at boot (bus-failure test, 22 Oct) & v0.7.7 re-initialises after 1\,s of failed or stuck reads (flashed 22 Oct; recovery confirmed on the bench on 22 Oct with v0.7.8: four wire pulls, re-initialised each time, no reset) \\ \hline
Raw stream filled 85\,\% of the front node's 115\,200\,baud line & Hex text at 36 bytes a sample & Console at 921\,600\,baud \\ \hline
Front sync pulse not seen on the Raspberry Pi~5 & No header pin toggled while the node drove GPIO15: wire not connected or on another pin (not settled) & Jumper moved to GPIO5 and firmware built for GPIO5; pin~11 toggles every 500\,ms \\ \hline
Commands from the Pi acted on 20 to 25\,s late & Read in the modem task, which waits on MQTT timeouts without network & v0.7.8 command task (flashed 22 Oct) \\ \hline
\end{longtable}

\section{Fault investigations}
\label{sec:investigations}

These are the eight faults of Table~\ref{tab:faults} that needed a diagnosis, because more than one cause could explain what was seen; Table~\ref{tab:investigations} summarises them. Each is told as symptom, candidate causes, deciding test, cause and fix, and consequence.

{\small
\begin{longtable}{|>{\raggedright\arraybackslash}p{2.1cm}|>{\raggedright\arraybackslash}p{3.6cm}|>{\raggedright\arraybackslash}p{4.3cm}|>{\raggedright\arraybackslash}p{3.7cm}|}
\caption{Fault investigations: candidate causes, the test that separated them and the result.}
\label{tab:investigations}\\
\hline
\textbf{Fault} & \textbf{Candidate cause} & \textbf{Test} & \textbf{Result} \\ \hline
\endfirsthead
\hline
\textbf{Fault} & \textbf{Candidate cause} & \textbf{Test} & \textbf{Result} \\ \hline
\endhead
Boot loop (\ref{sec:inv-flash}) & PSRAM setting & Add \texttt{qio\_opi} on top of the flash change & Same loop, same program counter: not the cause alone \\ \cline{2-4}
 & Application code or hardware & Remove every override & Runs, modem works: cleared \\ \cline{2-4}
 & Flash access mode & ROM banner; chip detection; 16\,MB with DIO & DIO reported; boots, LittleFS mounts: cause \\ \hline
GNSS (\ref{sec:inv-gnss}) & Reply layout not as assumed & Raw replies indoors and with a rooftop fix & 9 fields empty, 17 or 18 with a fix, decimal degrees: cause \\ \cline{2-4}
 & GNSS warm-up & Poll a modem already running for minutes & Polls answered: normal 60 to 90\,s delay \\ \cline{2-4}
 & Modem switched off & Five-minute capture after reprogramming & No reply to GNSS or MQTT commands: boot pulse toggled the modem off \\ \hline
BMS (\ref{sec:inv-bms}) & Pack absent or scanner idle & Log every advertiser & Nothing at all logged: scanner idle \\ \cline{2-4}
 & Characteristics under another service & Compare with a teammate's parser & Found directly there: cause of the drop \\ \cline{2-4}
 & Stuck BLE stack on the ESP32 & Full power cycle & Same pattern: ruled out \\ \cline{2-4}
 & Checksum formula & Firmware's own query frame & Register byte excluded; replies failed by 3: cause \\ \hline
Invalid JSON (\ref{sec:inv-json}) & Damage in transfer & Indexed, length-checked, repeated read-back & All 6,020 records intact: ruled out \\ \cline{2-4}
 & Records wrong as stored & Parse each record & 5,545 contain \texttt{nan}: cause \\ \hline
Colour cast (\ref{sec:inv-colour}) & Detector channel order & Still with the standard tuning & Also pink: ruled out \\ \cline{2-4}
 & No infrared-cut filter & Still with the NoIR tuning & Wall neutral: cause \\ \hline
ToF null, then hang (\ref{sec:inv-tof}) & One sensor not answering & I\textsuperscript{2}C bus scan & Only the left sensor present; start-up stopped at the next \\ \cline{2-4}
 & Driver waiting for ever & Driver's default I/O timeout & Zero, wait for ever: cause of the freeze \\ \hline
Pi~4 under-voltage (\ref{sec:inv-power}) & Pi supply under detector load & 2\,s flag log, nothing attached & Flags clear: ruled out \\ \cline{2-4}
 & ToF board on I\textsuperscript{2}C & Board disconnected throughout & Under-voltage still occurred: ruled out \\ \cline{2-4}
 & Telemetry unit on Pi USB & Plug in during the log & Under-voltage 14\,s later, recurring: cause \\ \cline{2-4}
 & & Unit on Pi USB, then moved off (10 Oct) & 3 restarts in 20 min, then none in 4\,h \\ \cline{2-4}
 & Heat & Same log & Soft limit from 17:45, before the unit: separate; fan fitted \\ \hline
Consumer stall (\ref{sec:inv-stall}) & Modem or BLE work starving the consumers & Read the code: where each task runs and how it waits & All on core~0 and every wait yields: ruled out \\ \cline{2-4}
 & Consumers stopped on core~1 & Which task writes \texttt{imu\_hz} & Only the statistics task: frozen at 200, so it had stopped; drop rates equal the IMU tasks' own rates: cause \\ \cline{2-4}
 & & Bus failure on v0.7.6 (22 Oct) & 17,080 read failures, 0 drops, consumers 2 to 29\,ms old: fixed \\ \hline
\end{longtable}
}

\subsection{Boot loop after enabling the full flash}
\label{sec:inv-flash}

\emph{Symptom.} The build tool treated the board as an 8\,MB part without PSRAM, although the module has 16\,MB of flash and 8\,MB of PSRAM. Setting the flash size to 16\,MB with a matching partition table made the processor reset about once a second, always at the same program counter.

\emph{Candidate causes and tests.} The PSRAM setting, the flash override, the code or the hardware. One setting was changed at a time. Adding a PSRAM type (\texttt{qio\_opi}) left the same loop at the same program counter, so PSRAM was not the cause on its own. Removing every override restored normal running, which cleared the code and the hardware. Two independent readings then pointed to the flash access mode: the upload tool's chip detection confirmed 8\,MB of PSRAM, and the ROM boot banner reported the flash in dual-I/O (DIO) mode, while the setting tried asked for quad-I/O.

\emph{Fix and consequence.} The firmware used about 9\% of the default partition, so the stock configuration was kept until local logging needed the 16\,MB table. Then 16\,MB with DIO mode, set together with the table, booted and mounted the logging partition, confirming the diagnosis. Size, mode and table must be set together.

\subsection{Modem power key, GNSS warm-up and reply layout}
\label{sec:inv-gnss}

\emph{Symptoms.} Three problems overlapped. Indoors every GNSS poll had the wrong number of fields; for 60 to 90\,s after the GNSS was switched on, polls got no reply; on a rooftop the first fix was reported near latitude and longitude zero (0.25 and 0.10 degrees); and a later five-minute capture got no reply to any of its 69 polls.

\emph{Candidate causes and tests.} A reply layout other than the 13 fields assumed, a normal receiver start-up delay, or a modem that was not running. Because the parser logged the raw reply, the real layout was visible: nine empty fields indoors and, with a fix, 17 fields (18 while settling) in decimal degrees, so the old mapping had read satellite counts as the position. A modem that had run for several minutes answered every poll, and during warm-up it answered other commands at once, so only the GNSS was waiting. The silent capture did not fit warm-up, as MQTT commands also went unanswered: the power key toggles the modem, which keeps its supply when only the processor restarts, so the pulse sent at every boot switched a running modem off after each reprogramming.

\emph{Fixes and consequence.} Since v0.5.3 the parser anchors on the latitude hemisphere letter (Section~\ref{sec:gnss-parsing-validation}); since v0.5.2 the key is pulsed only if the modem does not answer \texttt{AT}. Tests must allow the warm-up, during which a zero poll count does not mean that polling has stopped.

\subsection{Battery-management connection and checksum}
\label{sec:inv-bms}

\emph{Symptoms.} In turn: no BLE advertisements from any device; once the pack was powered, a link that dropped before any data; mostly failed connections; and on the vehicle three frames, all rejected on the checksum.

\emph{Tests and causes.} Every advertiser was logged, so seeing no device at all in a room full of them showed that the scanner was idle: scanning depended on an address set only after the pack had been found. A teammate's independently written parser finds the characteristics directly, which showed that in the ESP32 library they are not under the assumed service; every service is now searched. A power cycle of the board did not change the failed connections, which ruled out a stuck BLE stack. The checksum formula came from community reverse engineering. The query frame the firmware itself sends, \texttt{DD A5 03 00 FF FD 77}, carries $\mathtt{0x10000} - (\mathtt{0x03} + \mathtt{0x00}) = \mathtt{0xFFFD}$, which leaves out the byte in the register position; the firmware included it, so every genuine reply failed by exactly three.

\emph{Fix and consequence.} Firmware v0.5.4 corrects the sum and prints the bytes of any frame that still fails; 14 consecutive valid records followed on the vehicle. The protocol is confirmed, but the signal on the vehicle, $-88$ to $-93$\,dBm, is near the limit for a link, so the unit may have to be mounted near the pack.

\subsection{Invalid JSON found only by reading the log back}
\label{sec:inv-json}

\emph{Symptom.} At the first read-back of the stored log, on 1 October, 5,545 of 6,020 records failed to parse as JSON. Earlier transfer attempts had lost characters.

\emph{Candidate causes and test.} Damage in transfer, over a serial link that also carries other output, or records wrong as stored. The dump command prints each record with its index and length, and the laptop script (\texttt{BenchTest/log\_dump.py}) discards any record whose length does not match and repeats until every index has arrived intact. After four passes all 6,020 records were intact, so transfer was ruled out. Every failing record contained \texttt{nan}: a field with no value, such as the position before a fix, was printed as \texttt{nan}, which JSON does not allow.

\emph{Fix and consequence.} The same text had been published over MQTT, and watching the live output had not shown it. Since v0.6.3 such fields are written as \texttt{null}, and all 6,020 records parse. A strict receiver would have rejected 92\% of the records; later link tests count the lines that fail to parse at the receiving end (Appendix~\ref{app:evidence}).

\subsection{Camera colour cast}
\label{sec:inv-colour}

\emph{Symptom and candidate causes.} The live view looked reddish. A test was written for each of three causes: the detector's channel order, the camera itself, and the module's lack of an infrared-cut filter, which has its own tuning file.

\emph{Tests and cause.} A still with the standard tuning showed a white wall strongly pink, as did the detector's frames; with the tuning for a module without the filter the wall was neutral. The cause was the module, not the channel order.

\emph{Fix and consequence.} That tuning alone washed the colours out (a blue card at 120, 124 and 125 on red, green and blue); a saturation of 1.8 restored them (Table~\ref{tab:colour-cards}). Black objects then looked purple, since many black materials reflect infrared; greying near-black pixels took 84 to 171\,ms a frame, too slow for the detection loop, so it is done in the live view only. The colours must be checked again in daylight, where there is more infrared. Whether the cast lowered detection confidence was not measured.

\subsection{Time-of-flight start-up hang}
\label{sec:inv-tof}

\emph{Symptoms.} On 6 October every distance from the three sensors was null; after a first fix, on 10 October the reader froze.

\emph{Tests and causes.} A scan of the I\textsuperscript{2}C bus showed only the left sensor, at its new address, and nothing when the ahead sensor was woken. One sensor was not answering, but the start-up, which gives each sensor its address in turn, stopped at the first failure and lost all three. The reader was changed to skip and log such a sensor, and the bridge to look for it again every 60\,s. That relies on an error being raised, and on 10 October none was: the driver's I/O timeout defaults to zero, which means wait for ever, so one stuck sensor froze all three.

\emph{Fix and result.} Each sensor now has a 0.5\,s I/O timeout. With it, the reader no longer froze but logged all three sensors as timing out, while a device stayed at the shared start-up address 0x29 throughout, even with every shutdown line held low: one sensor was not under the control of its shutdown line. After the sensor board's wiring was re-seated, all three opened, and in a three-minute hand test the ahead and right sensors read 80 to 264 and 157 to 199\,mm with a hand about 20\,cm away (\texttt{BenchTest/\allowbreak{}logs/\allowbreak{}2026-10-10\_\allowbreak{}tof\_\allowbreak{}hand\_\allowbreak{}test.txt}).

\emph{Consequence.} The loose wire was inferred from these symptoms, not seen. If a shutdown (XSHUT) wire comes off, its pull-up keeps that sensor on at the shared start-up address, where it collides with each sensor being woken, and none of the three starts (Appendix~\ref{app:wiring}). The timeout turns this into a logged failure, visible before a test, but cannot keep the other two running.

\subsection{Raspberry Pi~4 under-voltage and reboots}
\label{sec:inv-power}

\emph{Symptoms.} Through much of the bench work the Raspberry Pi~4 reported under-voltage and throttled to 600\,MHz. On 4 October the detector ran at 6.1\,frames/s, against 14.6 later on a healthy Pi with the same services (\texttt{BenchTest/\allowbreak{}logs/\allowbreak{}2026-10-10\_\allowbreak{}bench\_\allowbreak{}observations.md}). On 6 October the Pi stopped booting until everything was disconnected.

\emph{Candidate causes.} The Pi's own supply under the detector's load; the time-of-flight board on its I\textsuperscript{2}C bus; the telemetry unit on its USB port; or heat, since it had no cooling.

\emph{Test.} A logger recorded every 2\,s the power and temperature flags, the I\textsuperscript{2}C addresses and the USB serial devices, so that each change in the flags could be matched to a device (\texttt{BenchTest/\allowbreak{}logs/\allowbreak{}2026-10-06\_\allowbreak{}pi4\_\allowbreak{}power\_\allowbreak{}watch.log}). With nothing attached and the detector running, the flags were clear, which cleared the supply. The time-of-flight board stayed disconnected. From 17:45 the soft temperature limit came and went (65 to 69\,\textdegree C) with no under-voltage. The telemetry unit was plugged in at 17:58:05; 14\,s later, when its modem came up, under-voltage and throttling began and recurred every few seconds, and the detector slowed from 14.3 to 12\,frames/s.

\emph{Cause.} The modem's transmit current, drawn through the Pi's USB port, pulls the Pi's supply down; the current was not measured. On 10 October, with the unit on the Pi's USB port, the Pi restarted three times in about 20 minutes; after the unit was moved to the laptop (through a small hub without its own supply), there were none in 4 hours and the flags stayed clear. The boot failure of 6 October was probably the same fault but was not reproduced. Heat was separate: 81\,\textdegree C without cooling, 48 to 66\,\textdegree C with a fan.

\emph{Fix and consequence.} The fault lay in how the parts were powered together, not in any one part. The telemetry unit needs its own supply, such as a powered hub, on the bench and on the vehicle. On a 3\,A supply the Raspberry Pi~5 limits its USB ports to 600\,mA in total (its supply on the vehicle is still to be confirmed), so the same applies there, and a measured power budget against the 72\,V source is still needed.

\emph{Recurrence in the closed housing (16 October).} With the Raspberry Pi~4 in the printed box, still powering the telemetry unit from its USB port, its temperature rose from 68.6 to 73.5 and then 76.9\,\textdegree C over about 70 minutes with the detector running. The flags showed under-voltage, throttling and the soft temperature limit since boot (0xD0000). Later the same day it restarted while under-voltage was present at that moment (0x50005): a brown-out. In the same period the share of unreadable dashboard lines rose from about 2 to 10\,\% once the battery-management link became active. Two changes followed. The box had vents only in its floor and low in one wall, so both pods' lids gained an exhaust grille over the computer under a push-on rain hood (Appendix~\ref{app:mechanical}); and the telemetry unit must have its own supply, as found above. The temperature in the vented box and the line-loss figure with a separate supply have not yet been measured. The Raspberry Pi~5 also restarted four times during the mounting work on 19 October. A power logger recording its supply, flags and temperature every 5\,s showed normal values in the last reading before each restart (4.95 to 5.08\,V, flags clean, 57 to 59\,\textdegree C). The supply's lowest reading in the two minutes before each was 4.89 to 4.97\,V. Those were abrupt losses of power, from a loose connector, the source or the supply's overload protection, not a brown-out or overheating.

\emph{Separate supply and a wired serial link (21 October).} During the road session the telemetry unit was still powered from the Raspberry Pi~4's USB port, which also carried its data. It was then given its own battery, and the Raspberry Pi~4 started with no under-voltage messages, against 11 to 16 per start before. USB could no longer be used for data, because the board charges its battery from that port and would load the Pi again. The data moved to three wires instead: the unit's spare pins IO41 and IO42 to the Pi's UART on header pins 10 and 8, and a common ground, with no supply wire between them. Firmware v0.7.4.1 sends every dashboard line on this link at 460\,800\,baud, drops and counts a line rather than wait if the link's buffer is full, and accepts the local commands from the Pi. On the Pi, Bluetooth, which nothing uses, was disabled so that the full UART rather than the clock-dependent mini UART is on those pins, and the serial console was removed from them. In the first check the bridge read 216 lines without a bad one, and 2,489 by midday, against 2 to 11\,\% lost on USB and a third of the per-second records missing on the road session (Section~\ref{sec:res-acq}); 27 records from that session also had damaged fields. Running on its own battery, however, the unit restarted by itself, after which its modem did not answer for the rest of the morning, which fits the battery sagging under the modem's transmit bursts. A battery alone is therefore not the answer: the unit needs a regulated supply from the vehicle.



\subsection{Consumer stall on the rear unit}
\label{sec:inv-stall}

\emph{Symptom.} On 21 October, about 23 minutes after the telemetry unit was put on its own battery with its record on the wired link, both IMUs' drop counts began rising, by about 167 and 46 a second, until a reset; after one reset IMU2 was not detected. On 22 October, with no USB host attached, the same fault returned at about 168 and 46 a second (733,399 and 198,524 dropped by 12:58). In both episodes \texttt{imu\_hz} stayed at exactly 200 for both sensors, the vibration value was null and the modem did not answer.

\emph{Candidate causes.} (a) Modem or BLE work taking processor time from the tasks that consume the samples, since the modem was silent both times; (b) a lock shared with those tasks, or the heap running out; (c) the USB console blocking, ruled out on 22 October because no host was attached; (d) the consumers themselves, the statistics task and the two ride tasks, having stopped.

\emph{Deciding evidence.} Reading the code of v0.7.4.1 decided it. \texttt{imu\_hz} is written only by the statistics task at the end of each 2\,s report, so a value frozen at 200.0 while the drops climbed meant that task had stopped; a null vibration value meant both ride tasks had stopped too. The statistics task drains both queues in one pass, so a slow but running consumer would drop the same number from both 200\,Hz producers; unequal rates cannot come from the consumer side. With the consumer stopped, every sample sent is dropped, so the drop rates were the IMU tasks' own rates: both were behind their 5\,ms schedule, IMU2 at about 22\,ms a pass. A task behind schedule never sleeps, and all three consumers ran on core~1 below the IMU tasks (priority~3), so they got no time. The modem, MQTT, GNSS and BLE work all run on core~0 and every wait there yields; no lock is shared with the consumers and nothing on the sample path allocates. That ruled out (a) and (b). What loaded core~1 was not proven. The candidates were interrupt load from the Pi's UART receive line floating while the Pi rebooted, and IMU2's I\textsuperscript{2}C bus. The same IMU2 later ran on the front node, on short wires, at 200\,Hz with no read failures, which points to its long wiring on the rear unit.

\emph{Fix.} Firmware v0.7.6 moved the statistics and ride tasks to core~0 at priority~2, above the modem, BLE and dashboard tasks, so that acquisition and processing no longer share a core. It also bounds the IMU catch-up at 20 periods, computes rates from the real elapsed time, adds the consumer ages, highest queue fill, late slots and idle counts per core to the dashboard line (Section~\ref{sec:imu-health}), and adds a command to switch the modem off for the bench test.

\emph{Field evidence and consequence.} On 22 October, with v0.7.6 on the rear unit, IMU1's wires were taken off while it ran and replaced with shorter ones: an unplanned but real bus failure. The unit counted 17,080 read failures and 8,508 late slots, but drops stayed at 0, the consumers were 2 to 29\,ms old and both cores kept idling, so the stall did not recur. After reconnection the sensor came back in its power-on sleep state and returned the same reading every time (stuck count rising 200 a second, vibration 0.0) until the board was reset, because the firmware configured a sensor only at boot. That led to the automatic re-initialisation of v0.7.7 (Section~\ref{sec:imu-health}), confirmed on the bench on 22 October with v0.7.8: in four pulls of IMU1's supply wire the sensor was re-initialised each time with no reset and no drops (Section~\ref{sec:impl-health}). The lessons are that processing must not share a core with acquisition and that a lost sensor must be recovered without a reset. The soak test of v0.7.6 is reported in Section~\ref{sec:res-acq}.

\subsection{GNSS polls starved during broker failures}
\label{sec:inv-gnss-hang}

\emph{Symptom.} On 24 October the rear node, off the vehicle with its antenna under open sky, held a 3D fix for 10 minutes (Section~\ref{sec:res-position}). Later that evening, while it was carried on a walk, every GNSS poll from 22:49:27 timed out (``poll timeout'' every 3\,s) instead of reporting a fix or no fix. Throughout, its MQTT reconnections to a broker that accepts only TLS failed (\texttt{CMQTTCONNECT} code 3) with backoff. After a reset over USB the modem answered again, but the receiver took about 50 minutes to report a fix, held it for about a minute and lost it.

\emph{Candidate causes.} (a) The single modem task blocked by the failing broker exchanges, so that GNSS queries go unanswered; (b) the modem's GNSS engine stopped; (c) for the slow fix after the reset, a cold start with the antenna's sky view partly blocked while it was carried.

\emph{Deciding evidence.} On 25 October, with the firmware pointed at the public broker, the connection still failed (\texttt{CMQTTCONNECT} codes 15 and 3, \texttt{CMQTTACCQ} errors, each attempt about 21\,s), and the assisted-GNSS download failed in the same way, while the same broker accepted the laptop in 0.4\,s; so the broker was not the cause, and the SIM's data connection is the likely one. The statistics lines showed the GNSS polled about once in 30\,s instead of every 3\,s while the reconnections ran, and the receiver fixed (14 satellites, HDOP 1.9) in a gap between two attempts: cause (a) is confirmed, the modem task's broker exchanges starve the GNSS polls.

\emph{Status and fix.} Planned: a command to switch publishing off, so that a failed data connection cannot starve the GNSS, and a check of the SIM's data. Until then the run-2 procedure powers the node outdoors at least 10 minutes before the first pass, with a fix confirmed.

\section{Diagnosis records}
\label{sec:diagnoses}

Two further faults are recorded here in more detail than Table~\ref{tab:faults} gives, because the method of finding them is part of the evidence.

\subsection*{Serial routing on the target board}

After the IMUs were moved onto the target board, nothing beyond the ROM boot banner appeared on the serial port, across reflashes, resets and reseated wiring. Three hypotheses were tested and ruled out in turn: a hung I\textsuperscript{2}C bus (a \texttt{Serial.flush()} immediately after \texttt{Serial.begin()}, before any I\textsuperscript{2}C code, still printed nothing); USB enumeration time (a fixed start-up delay changed nothing); and damage from the rewiring (power LED and temperature normal). The cause was a build setting made earlier for the bare development board, which reaches the host over UART0 through a USB-to-serial chip: on the target board the schematic shows UART0 wired through a CH340K to a separate USB-C connector, so application output went to the unconnected port while the ROM banner, which always uses native USB, reached the connected one. Restoring the native-USB setting fixed it at once; the two boards need opposite settings, which are now documented in \texttt{platformio.ini}.

\subsection*{Raspberry Pi~4 provisioning and network}

A boot drive made by copying boot files had one FAT32 partition and a \texttt{cmdline.txt} pointing to a root partition that did not exist; inspecting the partition table rather than the file listing showed it, and writing the image with Raspberry Pi Imager gave the expected boot (about 512\,MB) and root (about 29\,GB) partitions, checked against the disk signature. The Pi did not join a 5\,GHz laptop hotspot: \texttt{iw reg get} showed the kernel's country code but a generic ``world'' domain on the radio itself, and \texttt{rfkill} showed no block. Moving the hotspot to 2.4\,GHz and setting the country code through the driver-level configuration restored the link; both changes were applied together, so the sufficient fix was not isolated. A second provisioning pass, meant only to change the network credentials, had left ssh disabled, which shows that each field of a staged configuration has to be re-checked rather than assumed.

\section{Supporting detail}

\subsection*{GNSS reply layout}

With a fix, the fitted modem's \texttt{+CGNSSINFO} reply has seventeen fields: fix mode; GPS, GLONASS, BeiDou and Galileo satellite counts; latitude and hemisphere; longitude and hemisphere; date; UTC time; altitude; speed; course; and PDOP, HDOP and VDOP, with latitude and longitude in decimal degrees. It matches the sixteen-field form in version 1.09 of the manual \cite{simcoma76xx} with one extra satellite count; while a fix was settling an eighteen-field variant with an empty course and one undocumented trailing value also appeared, and the nine-field reply seen indoors is the same reply with the empty satellite block collapsed. The first outdoor capture, with a second query competing with the firmware's polling, contained replies with dropped characters, which the jump check now rejects. The modem's start-up command sequence takes about 27\,s after boot before polling begins (Section~\ref{sec:inv-gnss}).

\subsection*{Camera timing and modes}

The lens-cover test covered the lens by hand during a recording, with the operator reporting the instant and the Raspberry Pi's wall clock read at once; that time was compared with the recording's start time plus the elapsed time to the first frame that turned uniformly dark. A first attempt in which the operator reacted to a countdown failed (the event was missing in one case and seconds late in another); with the operator triggering and reporting the event, the two agreed to within one second. The OV5647 sensor (2592$\times$1944 pixels) supports, among other modes, 1920$\times$1080 at about 32.8 frames/s and 1296$\times$972 at about 46.3; the detector reads it at 640$\times$480 for speed.

\subsection*{Planned VL53L5CX validation}

When the VL53L5CX arrives it is to be validated on the bench before use: (1) confirm I\textsuperscript{2}C communication and the device identifier; (2) repeat the white and black target comparison of \cite{caroleo2026} at a fixed distance to measure this unit's valid-reading rate and bias; (3) sweep a target through the rated range to build a measured distance curve; (4) confirm that a target moved left, centre and right is reported in the right zones; and (5) check whether the warm-up drift affects the alert threshold. A single unit is planned because its 8$\times$8 zones give the left, centre and right discrimination; a second would be needed only for a wider field or another direction, such as reverse.

\subsection*{Raspberry Pi~4 discovery}

The dashboard's resolver (\texttt{DashboardIntegration/\allowbreak{}sw7\_\allowbreak{}endpoints.py}) tries, in order: localhost, an explicit override, a fixed address on the Ethernet link (10.20.0.1 and 10.20.0.2, in use since 15 October), the source of a UDP beacon sent by the bridge every 2\,s on port 50808 (ignored if older than 10\,s), \texttt{pi4-camera.local} through mDNS, and the last address that worked (Figure~\ref{fig:pi4-discovery}). A candidate is accepted when it answers with valid JSON within 3\,s; three failed polls start a new search, with the wait doubling to 30\,s. The offscreen tests pass; the wired link has been in use since 15 October and carried the rear data on the road session of 19 October.

\begin{figure}[htbp]
\centering
% How the dashboard finds the Raspberry Pi 4. Source: DashboardIntegration/sw7_endpoints.py
% (candidate order, PROBE_TIMEOUT_MS = 3000, FAILURES_BEFORE_REPROBE = 3, MAX_RETRY_MS = 30000),
% telemetry_bridge.py (beacon every 2 s on UDP 50808), live_data_provider.py (POLL_MS = 500).
\begin{tikzpicture}[x=1cm, y=1cm]
  \def\w{1.8cm}
  \node[swbox, text width=\w] (k1) at (1.25,0) {1. localhost\\ (one computer)};
  \node[swbox, text width=\w, right=0.32cm of k1] (k2) {2. override\\ (set by hand)};
  \node[swbox, text width=\w, right=0.32cm of k2] (k3) {3. wired link\\ 10.20.0.1};
  \node[swbox, text width=\w, right=0.32cm of k3] (k4) {4. beacon\\ UDP 50808};
  \node[swbox, text width=2.75cm, right=0.32cm of k4] (k5) {5. mDNS name\\ \texttt{pi4-camera.local}};
  \node[swbox, text width=\w, right=0.32cm of k5] (k6) {6. last address\\ that worked};
  \foreach \a/\b in {k1/k2, k2/k3, k3/k4, k4/k5, k5/k6} \draw[swflow] (\a) -- (\b);
  \begin{scope}[on background layer]
    \node[swlane, fit=(k1)(k6), inner sep=6pt] (search) {};
  \end{scope}
  \node[swlanetitle] at (search.north west) {Search, in this order; a candidate is accepted when it returns valid telemetry within 3\,s};
  \node[swbox, text width=4.6cm] (use) at (3.0,-2.6) {Connected: poll every 0.5\,s,\\ save the address as the last that worked};
  \node[swbox, text width=4.6cm] (wait) at (11.6,-2.6) {Nothing answered:\\ wait 2, 4, 8, 16, then 30\,s at most};
  \draw[swflow] (search.south -| use.north) -- node[swlabel, right=2pt] {first valid reply} (use.north);
  \draw[swflow] (k6.south |- search.south) -- node[swlabel, left=2pt] {none replied} (k6.south |- wait.north);
  \draw[swflow] (use.west) -- ++(-0.6,0) |- node[swlabel, pos=0.25, right=3pt, align=left] {3 failed polls\\ in a row, about 1.5\,s} (search.west);
  \draw[swflow] (wait.west) -- ++(-0.9,0) |- node[swlabel, pos=0.3, left=2pt] {search again} (search.south -| 9.0,0);
\end{tikzpicture}

\caption{How the dashboard finds the Raspberry Pi~4 (\texttt{DashboardIntegration/\allowbreak{}sw7\_\allowbreak{}endpoints.py}). The same resolver serves the camera page and the live-data provider.}
\label{fig:pi4-discovery}
\end{figure}

\subsection*{Camera colour}

With the NoIR tuning and saturation 1.8, printed cards held about 1\,m from the camera measured as in Table~\ref{tab:colour-cards} in the detector's own frames under indoor light. The three colours are clearly identified; grey and white show a small cool shift. Held together in one frame the cards remained distinct. The dark-pixel correction (Section~\ref{sec:inv-colour}) scales each pixel's saturation from fully grey at a brightest channel of 50 or below to unchanged at 110 and above, so the cards pass through it unchanged. The colours should be checked again in daylight.

\begin{table}[htbp]
\centering
\caption{Colour cards in the detector's frames with the NoIR tuning and saturation 1.8 (mean red, green and blue, 0 to 255; 1 October 2026, indoor light).}
\label{tab:colour-cards}
\begin{tabular}{lrrr}
\hline
Card & R & G & B \\
\hline
Blue & 1 & 86 & 216 \\
Red & 222 & 8 & 77 \\
Green & 6 & 124 & 103 \\
Grey & 124 & 129 & 138 \\
White & 186 & 199 & 197 \\
\hline
\end{tabular}
\end{table}

\subsection*{Pothole model}

Run 1 (2 October 2026) fine-tuned YOLO11n, pretrained on COCO, on a public 608-image pothole dataset from Roboflow Universe \cite{gerapothole} in Colab, with no images from this vehicle's camera; its test images come from the same source, so its score describes that dataset, not South African roads. It ran 46 epochs in 15 minutes, stopped 15 epochs after the best validation result; mAP@50 0.597 at 640 pixels (precision 0.66, recall 0.50) and 0.550 at 320 on the dataset's own test images. Validation loss stopped improving after about 20 epochs, so more varied images were needed, not longer training. On the South African dataset's 605 test images (1,238 potholes) it scored 0.000 at 640, 1280 and 1920 pixels, recall at most 1.2\%; drawing the converted labels on an image confirmed that the label conversion was correct. Of that dataset's 4,337 images (3680$\times$2760 pixels), 1,987 contain 5,029 potholes, with a median size of 100 by 26 pixels, about 17 by 4.5 pixels after reduction to a 640-pixel input, and all centres lie between 48 and 62\% of the image height. Run 2 (6 October) cropped each frame to the band from 42 to 68\% of its height and cut it into three overlapping 1280-pixel tiles, so that a median pothole is about 50 by 13 pixels at the training input; after 60 epochs on a T4 it reached a validation mAP@50 of 0.714 at 640 pixels (precision 0.81, recall 0.62) and 0.566 at 320. These are validation scores, optimistic because validation chose the best epoch, and per tile; the notebook's record also did not state which weights training started from (both corrected in the notebook on 6 October). On the Raspberry Pi~4, installing the PyTorch library used by the Ultralytics software failed over the available network, so the exported model was run with the NCNN library alone, with letterboxing, box decoding and non-maximum suppression written for the purpose (\texttt{CameraDetection/\allowbreak{}ncnn\_\allowbreak{}bench.py}), the lighter form a deployment would take. In a 15\,s trial a COCO YOLO11n at a 320-pixel input with four threads ran at 7.9\,frames/s, with a median of 111\,ms for inference and 175\,ms from readout to result (95th percentile 258\,ms), against 19.5\,frames/s and 102\,ms for the SSD-MobileNet detector at its 300-pixel input.

\subsection*{Dashboard map and navigation}

The street-tile map (\texttt{DashboardIntegration/\allowbreak{}sw7\_\allowbreak{}tile\_\allowbreak{}map.py}) draws OpenStreetMap tiles with Qt's painter under the tile usage policy (at most two requests at a time, an identifying user agent, no bulk download), caches tiles for 30 days, falls back to scaled lower-zoom tiles or the team's vector roads, and follows heading-up at zoom 17. A prefetch tool cached 656 tiles (11.9\,MB) at zoom 11 to 16 around the campus, since the tile policy forbids bulk download at zoom 17 and above. Routes come from the public OSRM service \cite{luxen2011osrm}; map data are \copyright\ OpenStreetMap contributors \cite{osm}. Firmware v0.7.1 added the GNSS course over ground (\texttt{crs}) and altitude (\texttt{alt}) to the dashboard line for the heading-up map, both null without a fix; a recorded bench check of the v0.7.1 output is not in the repository. A generated replay track (2,993\,m along campus roads from an OSRM route, 650 points at one per second at 18\,km/h, not recorded data) and a simulated drive at 25\,km/h allow the map to be shown indoors; both are labelled on screen (REPLAY, Simulated drive) so that they cannot be taken for real positions, nothing from them is written to the vehicle state, and a real fix stops the simulation. The map has not yet followed a real fix. For bench work the start script also skips the driver and PIN page; the page is to be restored for the demonstration. The Raspberry Pi~5 also carries a Hailo-8 accelerator, which did not appear on the PCI Express bus because the external connector is not enabled in its boot configuration.
